\chapter{研究报告纠错核查}
\chapterlead{对应子选题“研究报告纠错核查”。判一个正确的数为错，比漏掉一个错数更伤信任。}

\section{要解决什么问题}

纠错核查的输入是一份研报和对应公司的年报，输出是一张逐条核查表：研报中每个可核验的数字断言，对还是错；错的话属于哪类错误，正确值是多少，
出自年报哪一页，应该怎么改。它服务的是研究所质控人员，所以有两个硬要求：误报要少，否则核查表会被当成噪声；
结论要具体，“该数字可能有误”这样的提示没有用，要说清楚是单位写错了、期间串了，还是把净利润当成了归母净利润。

\section{核查流程}

\begin{figure}[htbp]
\centering
\begin{tikzpicture}[node distance=4mm and 6mm]
  \node[doc,minimum width=1.7cm,minimum height=1.0cm,font=\small] (rep) {研报};
  \node[agent,right=of rep,yshift=6mm,minimum width=2.3cm] (llm) {模型拆断言\\ \scriptsize claims.json};
  \node[prog,right=of rep,yshift=-6mm,minimum width=2.3cm] (rule) {规则抽取\\ \scriptsize 中文数字模式};
  \node[prog,right=8mm of llm,yshift=-6mm,minimum width=2.0cm] (cl) {断言集\\ \scriptsize 指标·期间·值·单位};
  \node[prog,right=of cl,minimum width=2.2cm] (ver) {逐条核对\\ \scriptsize 容差 $=\halfulp$};
  \node[gate,right=of ver] (ok) {一致?};
  \node[prog,below=8mm of ok,minimum width=2.4cm] (diag) {错误类型诊断\\ \scriptsize 单位→方向→期间→口径};
  \node[prog,fill=progc!18,above=8mm of ver,minimum width=1.6cm] (tt) {真值表};
  \node[doc,right=8mm of ok,text width=2.0cm,font=\scriptsize] (out) {核查表\\ 结论 · 正确值 · 出处 · 修改建议};
  \draw[arr] (rep) -- (llm); \draw[arr] (rep) -- (rule); \draw[arr] (llm) -- (cl); \draw[arr] (rule) -- (cl);
  \draw[arr] (cl) -- (ver); \draw[arr] (tt) -- (ver); \draw[arr] (ver) -- (ok); \draw[arr] (ok) -- node[above,font=\scriptsize]{是} (out);
  \draw[arr] (ok) -- node[right,font=\scriptsize]{否} (diag); \draw[arr] (diag) -| (out);
\end{tikzpicture}
\caption{纠错核查流程。断言抽取可由模型完成，核对与诊断全部由程序完成。}
\label{fig:rc-flow}
\end{figure}

\subsection{断言抽取}

断言是研报里一句可以核对的数字陈述，拆成指标、期间、数值、单位和类型（金额、同比、比率、百分点变化）四五个要素，并带上原文和位置。
抽取有两条通道：规则通道用中文数字模式匹配“同比增长 25.30\%”“实现营业收入 58.62 亿元”这类句式；模型通道由分析师会话逐句拆出断言，
处理指标别名、“上半年”“报告期”这类期间指代和省略的单位。含“预计”“预测”“目标”或年份带 E 的句子标为预测性断言。

\subsection{按展示精度核对}

核对的关键是容差。研报写“58.62 亿元”，年报写“5,862,340,125.36 元”，两者不相等，但写研报的人只对百万元级别负责。
系统按研报数字的展示精度定容差：两位小数的“亿元”，容差是 0.005 亿元，即 50 万元；一位小数是 500 万元；百分数写两位小数，容差是 0.005 个百分点。
这个容差恰好是“这样写的人不可能知道更多”的范围。在容差内一致判“正确”；预测性断言和真值表里没有对应值的断言判“无法核实”，不计为错误。

\subsection{错误类型诊断}

数字对不上时，程序按固定顺序尝试替代解释，第一个精确命中的作为结论：
\begin{enumerate}
  \item \textbf{单位错误}：换成万元、元、百万元后与真值一致；
  \item \textbf{方向错误}：同比取相反数后与真值一致（“增长”写成了“下降”）；
  \item \textbf{期间错配}：与上一期或其他期间的真值一致；
  \item \textbf{口径错误}：与同一口径组里另一个指标的真值一致，如净利润、归母净利润、扣非归母净利润、利润总额、营业利润；
  \item 都不命中，判\textbf{数值错误}。
\end{enumerate}
每种替代解释也用同样的展示精度容差，必须精确命中才下结论。系统宁可只说“数值错误”，也不给一个猜测的错误类型。
修改建议随类型生成，例如口径错误会写“若指归属于上市公司股东的净利润，应为 6.12 亿元；若指净利润，请更正指标名称”。

\begin{story}[一次诊断的完整过程：样例研报 E1]
研报写“归母净利润 6.46 亿元”。真值表中 2024 年归母净利润为 612,380,000 元，差 3,380 万元，远超 50 万元的容差，不正确。
按顺序诊断：换单位，不命中；不是同比，跳过方向；2023 年值 518,210,000 元，不命中；口径组里“净利润”为 646,102,000 元，与 6.46 亿元相差 10.2 万元，
在容差内，命中，判“口径错误”。出处给到年报第 6 页的表块与行列。整个过程没有模型参与，每一步都可以复核。
\end{story}

\section{结果}

\subsection{样例研报}

样例公司“星河智造”是虚构的，年报与研报草稿由脚本生成，内部完全自洽，用于回归测试。草稿里埋了 7 处错误，系统全部判出，错误类型全部正确，误报 0；
另有 1 条预测性断言判“无法核实”（表~\ref{tab:rc-sample}）。整个评测不调用模型，耗时约 0.01 秒，一条命令即可复现。

\begin{table}[htbp]
\centering
\caption{样例研报的 7 处埋错与系统结论}
\label{tab:rc-sample}
\small
\begin{tabularx}{\linewidth}{@{}c l l l L@{}}
\toprule
\hd{编号} & \hd{研报原文} & \hd{金标准} & \hd{系统结论} & \hd{真值与依据} \\
\midrule
E1 & 归母净利润 6.46 亿元 & 口径错误 & 口径错误 & 6.12 亿元；6.46 亿元与“净利润”一致 \\
\rowband
E2 & 同比增长 25.30\% & 数值错误 & 数值错误 & 18.17\%，无替代解释 \\
E3 & 经营现金流 2.38 万元 & 单位错误 & 单位错误 & 2.38 亿元 \\
\rowband
E4 & 毛利率 32.50\% & 数值错误 & 数值错误 & 31.19\%（由营业收入与营业成本派生） \\
E5 & 提升 1.2 个百分点 & 数值错误 & 数值错误 & $-0.03$ 个百分点；方向也反，但取反后仍不符 \\
\rowband
E6 & 营业成本 35.22 亿元 & 期间错配 & 期间错配 & 40.34 亿元；35.22 亿元为 2023 年值 \\
E7 & 应收账款同比 28.50\% & 数值错误 & 数值错误 & 41.60\% \\
\bottomrule
\end{tabularx}
\end{table}

E5 值得多说一句。研报写“提升 1.2 个百分点”，真值是下降 0.03 个百分点，方向确实相反，但取反之后 $-1.2$ 与 $-0.03$ 仍然差得很远，
所以诊断落到“数值错误”。如果系统在这里贸然判“方向错误”，修改建议就会变成“把提升改成下降”，改完仍然是错的。

\subsection{在真实真值表上的注入实验}

7 处埋错太少，不足以比较策略。在人工埋错的真实研报评测集完成之前，我们在茅台、宁德时代、招商轮船和样例公司四张真值表上，用程序生成了一批断言：
正确断言按真值四舍五入写成一位或两位小数；含错断言分五类，数值错误按 $\pm0.3\%$ 到 $\pm30\%$ 随机偏移，单位错误把“亿元”写成“万元”，
期间错配写上年值，口径错误写同组另一指标的值，方向错误把同比取反。随机种子固定，共 1,334 条，其中正确 580 条、含错 754 条。

我们把系统的容差策略与五种固定容差对比（表~\ref{tab:rc-tol}、图~\ref{fig:rc-tol}）。按展示精度定容差，召回 100\%、误报 0；固定容差没有一个点能同时做到。
精确相等会把 580 条正确断言中的 563 条判错，因为舍入差；相对 0.1\% 仍有 28.4\% 的误报，因为“1.7 亿元”与真值 1.684 亿元相差 0.95\%，
小金额、少位数时舍入误差的相对值很大；放宽到相对 5\%，误报降下来了，但漏掉了 106 个小幅数值错误和 32 个期间错配。

\begin{table}[htbp]
\centering
\caption{容差策略对比（注入集：含错 754 条，正确 580 条）}
\label{tab:rc-tol}
\small
\begin{tabular}{@{}lrrrrrr@{}}
\toprule
\hd{容差策略} & \hd{判出} & \hd{误报} & \hd{召回率} & \hd{精确率} & \hd{误报率} & \hd{类型准确率} \\
\midrule
\textbf{展示精度 $\halfulp$（系统）} & 754 & 0 & \textbf{100.0\%} & \textbf{100.0\%} & \textbf{0.0\%} & \textbf{99.5\%} \\
\rowband
精确相等 & 754 & 563 & 100.0\% & 57.3\% & 97.1\% & 39.4\% \\
相对 0.1\% & 754 & 165 & 100.0\% & 82.0\% & 28.4\% & 95.2\% \\
\rowband
相对 1\% & 700 & 19 & 92.8\% & 97.4\% & 3.3\% & 97.7\% \\
相对 5\% & 608 & 3 & 80.6\% & 99.5\% & 0.5\% & 92.8\% \\
\rowband
绝对 0.5 & 663 & 0 & 87.9\% & 100.0\% & 0.0\% & 96.5\% \\
\bottomrule
\end{tabular}
\end{table}

\begin{figure}[htbp]
\centering
\begin{tikzpicture}
\begin{axis}[width=0.74\linewidth, height=5.4cm, xlabel={误报率（\%）}, ylabel={召回率（\%）},
  xmin=-5, xmax=105, ymin=76, ymax=103, grid=major, grid style={rule!50},
  nodes near coords, point meta=explicit symbolic,
  every node near coord/.append style={font=\scriptsize,anchor=west,xshift=2pt}]
\addplot[only marks,mark=*,mark size=3pt,color=seal] coordinates {(0.0,100.0) [展示精度（系统）]};
\addplot[only marks,mark=square*,mark size=2.3pt,color=inkl] coordinates {
  (97.1,100.0) [精确]  (28.4,100.0) [相对 0.1\%]  (3.3,92.8) [相对 1\%]  (0.5,80.6) [相对 5\%]  (0.0,87.9) [绝对 0.5]};
\end{axis}
\end{tikzpicture}
\caption{固定容差沿一条曲线移动：收紧则误报激增，放宽则漏掉小幅错误。按展示精度定容差位于左上角。}
\label{fig:rc-tol}
\end{figure}

容差还决定了“能不能说清错在哪里”。精确相等策略下，单位、期间、方向、口径四类错误都判出来了，但类型几乎全错（单位错误只有 6/132 判对），
因为替代解释也要求精确命中，舍入误差让它们全部失效，诊断落到了兜底的“数值错误”。系统策略下仅有的 4 个类型错误，
是注入在折旧摊销合计上的口径错误被判成了“数值错误”：这个派生指标当时不在指标字典里，诊断时被跳过。我们随后把它加入字典并补了回归测试。
系统没有给出错误的类型，只是没有给出更具体的类型。

\begin{note}[这个实验高估了真实效果]
注入集的正确断言恰好按真值四舍五入生成，含错断言只有一种错误，而且直接以结构化形式进入核对器，没有经过断言抽取。
真实研报还有三类难点：断言抽取本身的错误（指标别名、期间指代、单位省略）；作者的取整习惯与年报不同；一个数同时满足两种替代解释。
所以 100\% 与 0\% 应理解为“核对器在理想输入上不引入误差”，端到端的效果要等人工埋错的真实研报评测集完成后测量（第 13 章）。
\end{note}

\section{研报内部一致性}

除了与年报比对，研报自身也可能前后不一：摘要写“买入”，正文写“增持”；首页的 2026 年预测利润与估值章节的不一样。
内部一致性检查比较全文各处的评级、各情景每股价值、预测年度的公司整体收入与利润，只有实体、口径、期间、单位、情景都相同时才算矛盾。
引述外部机构观点的句子（带来源编号）不参与比较。这个检查同时用在三个地方：纠错核查、数字护栏，以及质检评分中的“数字前后矛盾”条目。
